\section{Background}

\subsection{Imputation methods in survey practice}

Survey agencies and microsimulation teams draw on a small set of imputation
families. \emph{Hot-deck and statistical matching} methods donate observed
values from a matched donor record; they preserve marginal distributions by
construction because every imputed value is a real donor value
\citep{andridge2010review, dorazio2006statistical}, and remain the standard
data-fusion tool in European microsimulation practice
\citep{sutherland2013euromod}. \emph{Regression-based} methods impute from a
fitted conditional model --- ordinary least squares with residual draws
\citep{vonhippel2007should}, linear quantile regression
\citep{koenker1978regression, koenker2005quantile} --- and extend to multiple
imputation for uncertainty propagation \citep{rubin1987multiple,
kennickell1998multiple}. \emph{Tree-ensemble} methods, in particular quantile
regression forests \citep{meinshausen2006quantile, breiman2001random},
estimate full conditional distributions nonparametrically and have become the
default in \policyengine{}'s data pipelines \citep{ghenis2018quantile,
woodruff2024enhancing} for their ability to capture nonlinear structure without
per-variable specification. The \microimpute{} package
\citep{policyengine2025microimpute} implements these families under one
interface and selects among them by cross-validation; its central finding ---
no family dominates across datasets --- motivates empirical comparison on
each new task rather than a prescribed method.

\subsection{Chained imputation and joint structure}

Imputing several variables one at a time, each conditional only on shared
predictors, destroys the dependence among imputed variables; sequential
(chained-equations) imputation conditions each variable on those already
imputed, retaining joint structure \citep{raghunathan2003multiple,
vanbuuren2011mice}. In data fusion the same issue appears as the conditional
independence assumption of statistical matching: matched files preserve
marginals but attenuate cross-source correlations \citep{dorazio2006statistical,
meinfelder2011simulation}. Synthetic-data generation faces the identical
trade-off at whole-file scale \citep{rubin1993statistical}.

\subsection{Weights in imputation}

Design weights enter imputation practice unevenly. The model-based literature
often treats weights as a nuisance --- if the model is correctly specified and
the design variables are conditioned on, weighting is unnecessary --- while
the design-based tradition weights every estimator \citep{little2002statistical}.
In production microdata pipelines the model is never correctly specified and
the design variables are never fully available, so the choice is consequential:
an unweighted fit targets the sample conditional, and any correlation between
weights and outcomes within predictor cells propagates into the imputed
population distribution. Tree ensembles add a mechanical subtlety: a
fully-grown forest's \emph{predictive distribution} is a leaf-membership
distribution, which per-record \texttt{sample\_weight} arguments do not
reweight, so honoring weights requires materializing them into the training
data --- the weighted bootstrap of Section~3. Zero-inflated targets add a
second subtlety: weighting the zero/nonzero gate by resampling can delete a
rare sign class outright, so the gate must be weighted directly. We are not
aware of prior survey-imputation work that combines weight-aware forest
training, sign-regime gating, and chaining in one estimator; documenting and
testing that combination is the gap this paper fills.

\subsection{Evaluating imputed files}

Because an imputation is a draw from an estimated conditional distribution,
pointwise accuracy metrics understate quality; the relevant question is
distributional. We follow the quantile-loss tradition for conditional
calibration \citep{koenker2005quantile, meinshausen2006quantile}, use
Wasserstein distance for marginal fit, and adopt the precision--density--coverage
framework of \citet{naeem2020reliable} for joint structure across imputed
columns, as prior \policyengine{} evaluation work did. The reweighting
fragility diagnostic of Section~3 is, to our knowledge, new.
